\section{Related Work}

Yan et al.\ combine a low-speed tractor--trailer kinematic model, reference
trajectory errors, preview states, and PPO for reverse path following
\cite{yan2024trajectory}. Their parameterization includes tractor wheelbase,
rear-axle-to-hitch offset, and hitch-to-trailer-axle distance. Their work
motivates the implemented trajectory-relative observation in this project, but its
software interface and several reproduction details are not adopted verbatim.

Moradi and Mehrandezh formulate an NMPC local planner for forward and reverse
tractor--trailer motion, including state and control constraints, articulation
limits, obstacle avoidance, and a transition from a headless Python model to a
Gazebo/ROS implementation \cite{moradi2022nmpc}. Their compact experimental
dimensions and staged simulator workflow motivate the nominal vehicle scale
and development sequence used here. Our primary comparison differs by using a
shared fixed-speed, steering-rate plant interface for PPO and NMPC.

Ma et al.\ separate long-horizon reference-line graph search from constrained
reinforcement-learning trajectory generation \cite{ma2025hierarchical}. This
supports a modular separation between path generation and local execution, but
multi-lane yard routing is outside the critical path of the present study.
Menz and Wenzel study PPO-generated demonstrations, behavioral cloning, and
offline reinforcement learning for a steerable drawbar-trailer configuration
\cite{menz2026imitation}. Their experiment-management and robustness questions
are relevant, although their vehicle mechanism differs from the conventional
semitrailer used here.

PPO provides a continuous-action policy-gradient baseline with clipped policy
updates \cite{schulman2017ppo}. NMPC will be implemented using CasADi
\cite{andersson2019casadi} and an interior-point nonlinear-programming solver.
Stable-Baselines3 is the planned PPO implementation \cite{raffin2021sb3}, and
MuJoCo is the implemented rigid-body physics backend \cite{todorov2012mujoco}.
